\documentclass[12pt]{article}

\usepackage[utf8]{inputenc}
\usepackage[T1]{fontenc}
\usepackage{amsmath, amssymb, amsthm}
\usepackage{geometry}
\geometry{a4paper, margin=2.5cm}

\title{Teoreme Fundamentale din Analiza Complexă}
\author{}
\date{}

\theoremstyle{plain}
\newtheorem{theorem}{Teoremă}

\begin{document}

\maketitle

\begin{theorem}
\textbf{Formula integrală a lui Cauchy.}  
Fie $f$ o funcție olomorfă în interiorul și pe conturul unei curbe închise netede $\gamma$, orientată pozitiv.  
Pentru orice punct $z$ din interiorul lui $\gamma$ are loc



\[
f(z) = \frac{1}{2\pi i} \int_{\gamma} \frac{f(w)}{w - z} \, dw.
\]



Aceasta exprimă valoarea funcției în interiorul conturului prin valorile sale pe contur.
\end{theorem}

\begin{theorem}
\textbf{Formula lui Cauchy pentru derivate.}  
În ipotezele teoremei precedente, pentru orice $n \ge 1$ și orice punct $z$ din interiorul lui $\gamma$ are loc



\[
f^{(n)}(z) = \frac{n!}{2\pi i} \int_{\gamma} \frac{f(w)}{(w - z)^{n+1}} \, dw.
\]



Aceasta arată că toate derivatele unei funcții olomorfe sunt determinate de valorile funcției pe contur.
\end{theorem}

\begin{theorem}
\textbf{Inegalitatea lui Cauchy.}  
Fie $f$ olomorfă în discul $D(z_0, R)$ și fie



\[
M = \max_{|z - z_0| = R} |f(z)|.
\]



Atunci, pentru orice $n \ge 0$ și orice $z$ cu $|z - z_0| < R$ are loc



\[
|f^{(n)}(z_0)| \le \frac{n! \, M}{R^n}.
\]



Aceasta oferă o estimare a derivatelor unei funcții olomorfe în funcție de valorile sale pe cercul de rază $R$.
\end{theorem}

\end{document}
